\begin{claim}{D.76}{Proved}{The full section augmentation module and its incidence map}
Retain the actual pants extension $E=F\rtimes_\alpha\langle s\rangle$
of D.74, with $F=\widehat F(x,y)$, and put
$k=\mathbb F_2$, $R=k[[E]]$, and $R_F=k[[F]]$.
Choose one actual subgroup $T_i=\langle t_i\rangle$ in each
of its three section classes, $0\leq i<3$.
The compact left $R$-module
\[
 \Delta_T=\ker\left(
       \bigoplus_{i=0}^2 k[[E/T_i]]
                    \xrightarrow{\ \varepsilon\ }k\right)
                                                               \tag{D.76.1}
\]
is projective and has a full continuous module isomorphism
\[
                         R^2\xrightarrow{\ \cong\ }\Delta_T .
                                                               \tag{D.76.2}
\]
The proof constructs this isomorphism from compatible
isomorphisms at every finite group-ring quotient.
For every finite continuous $E$-module $M$ over $k$,
\[
             \dim_k\operatorname{Hom}_E(\Delta_T,M)
                                      =2\dim_k M.             \tag{D.76.3}
\]

There is a separate, specified incidence map.
Let $e_i$ denote the identity coset in the $i$th summand
of (D.76.1), and put
$H=\overline{\langle T_0,T_1,T_2\rangle}$.
Then
\[
 \begin{aligned}
 d_T:R^2&\longrightarrow\Delta_T,
       &(a,b)&\longmapsto a(e_1-e_0)+b(e_2-e_0),\\
 \operatorname{coker}d_T
       &\xrightarrow{\ \cong\ }
       \ker\bigl(k[[E/H]]\xrightarrow{\ \varepsilon\ }k\bigr).
 \end{aligned}                                                 \tag{D.76.4}
\]
Under any isomorphism (D.76.2), $d_T$ is a square
module endomorphism. It is invertible exactly when
these three actual involutions generate $E$.
This criterion retains the chosen subgroup maps;
the existence of (D.76.2) does not assert invertibility
of the specified $d_T$.

In the normalized D.74 coordinates, four actual
involutions already generate $E$:
\[
                        s,\quad xs,\quad fs,\quad yfs.
                                                               \tag{D.76.5}
\]
The first and third belong to the same section class,
with the exact unique transporter
$fs=h^{-1}sh$ from D.74.
For the three actual involutions
\[
                   s,\qquad xs,\qquad hyh^{-1}s,             \tag{D.76.6}
\]
the subgroup $H\cap F$ in (D.76.4) is exactly
$\overline{\langle x,hyh^{-1}\rangle}$.
Generation by (D.76.6), and D.74's remaining third-cusp
equation, are not asserted by this module calculation.
\end{claim}

\begin{proof}
\emph{All actual coset and restriction maps.}
Every $T_i$ maps isomorphically to $E/F=C_2$, so
$E=FT_i$ and $F\cap T_i=1$. Hence
\[
        F\longrightarrow E/T_i,\qquad a\longmapsto aT_i
                                                               \tag{D.76.7}
\]
is an $F$-equivariant homeomorphism: it is a continuous
bijection between compact Hausdorff spaces.
It identifies the restriction of the $i$th permutation
module with the actual regular module $R_F$.
If $T_i=a_iS_i a_i^{-1}$ for a representative in D.74,
the full $E$-equivariant coset comparison is
$gT_i\mapsto ga_iS_i$, with inverse
$gS_i\mapsto ga_i^{-1}T_i$.
Thus all changes of the three representatives retain
explicit maps.

Write $I_F=\ker(R_F\to k)$.
In the coordinates (D.76.7), the mutually inverse
$R_F$-linear maps
\[
 \begin{aligned}
 I_F\oplus R_F\oplus R_F&\longrightarrow\Delta_T,
                &(a,b,c)&\longmapsto(a+b+c,b,c),\\
 \Delta_T&\longrightarrow I_F\oplus R_F\oplus R_F,
                &(a_0,a_1,a_2)&\longmapsto
                             (a_0+a_1+a_2,a_1,a_2)
 \end{aligned}                                                \tag{D.76.8}
\]
give the actual restricted module.
The augmentation ideal of the full free group is
free on $x-1,y-1$:
\[
       R_F^2\longrightarrow I_F,\qquad
                  (a,b)\longmapsto a(x-1)+b(y-1).
                                                               \tag{D.76.9}
\]
Here is a direct verification in the completed category.
For every finite $F$-module $M$, a continuous linear
map $I_F\to M$ is exactly a continuous crossed
homomorphism $F\to M$, by $g\mapsto g-1$ and
$d(gh)=d(g)+g\,d(h)$.
The images of $x,y$ of such a crossed homomorphism
are arbitrary elements of $M$: use the free profinite
universal property on the two lifts in the finite
semidirect product $M\rtimes\operatorname{im}(F\to\operatorname{Aut}M)$.
This also proves uniqueness.
Consequently (D.76.9) induces a bijection on continuous
Hom into every finite module. Applying those bijections
to the finite quotients of its source supplies compatible
inverse maps; finite quotients separate points.
Thus (D.76.9) is the asserted full continuous isomorphism.
In particular
\[
                         \Delta_T|_F\cong R_F^4.             \tag{D.76.10}
\]
It follows that $\Delta_T$ is finitely generated over
$R$ as well. No finiteness statement for a hypothetical
discrete group has been used.

\emph{A direct degree-one bridge for every coefficient module.}
Let $u:E\to C_2=k$ be the original quotient character,
and let $M$ be any finite continuous left $E$-module
over $k$. Give
$W_M=M\otimes_k k[E/F]$ the diagonal action.
The exact coefficient sequence is
\[
 \begin{gathered}
 0\longrightarrow M\xrightarrow{\ j\ }W_M
                    \xrightarrow{\ \epsilon\ }M
                       \longrightarrow0,\\
 j(m)=m\otimes(F+sF),\qquad
 \epsilon(m\otimes gF)=m.
 \end{gathered}                                               \tag{D.76.11}
\]
The actual induction comparison and its inverse are
\[
 \begin{aligned}
 R\widehat\otimes_{R_F}M&\longrightarrow W_M,
              &g\otimes m&\longmapsto gm\otimes gF,\\
 W_M&\longrightarrow R\widehat\otimes_{R_F}M,
              &m\otimes gF&\longmapsto g\otimes g^{-1}m .
 \end{aligned}                                                \tag{D.76.12}
\]
Changing $g$ to $gf$ with $f\in F$ leaves the second
formula unchanged by the tensor relation. Both composites
are the identity. Since the subgroup has finite index,
induction is also coinduction. Explicitly, $g\otimes m$
corresponds to the function on $E$ whose value at $a$ is
$(ag)m$ when $ag\in F$, and zero otherwise; the coinduced
action is right translation of the argument.
Shapiro therefore gives
$H^2(E,W_M)=H^2(F,M)=0$.

The connecting map from (D.76.11) is precisely
cup product by the named $u$. To check it without a
choice of quotient generator, lift a $1$-cocycle
$a:E\to M$ by $\widetilde a(g)=a(g)\otimes F$.
Its coboundary at $(g,h)$ is
\[
 j\bigl(u(g)\,g\,a(h)\bigr).
 \qquad\text{Thus}\qquad
 \partial[a]=[u\smile a].
                                                               \tag{D.76.13}
\]
The coefficient long exact sequence proves that
$u\smile-:H^1(E,M)\to H^2(E,M)$ is onto.

D.74 gives the actual equalities $N_E(T_i)=T_i$.
The finite-coefficient Brown theorem applies in degree
two: the open group $F$ has $\operatorname{cd}_2(F)=1$,
all finite subgroups inject into $E/F$, and the three
$T_i$ are exactly its full section classes.
The two-row Hochschild--Serre sequence with this finite
module makes $H^2(E,M)$ finite, so the dense embedding of
\cite[Theorem~2 and Section~1.1]{A-KochloukovaZalesskii2007}
is onto. Hence
\[
       H^2(E,M)\xrightarrow{\ \cong\ }
                            \bigoplus_{i=0}^2H^2(T_i,M).
                                                               \tag{D.76.14}
\]
On every actual $T_i$, $u$ is its nonzero quotient
character. The cyclic complex has differential $1+t_i$
in each positive degree. Cup product by $u$ identifies
$H^1(T_i,M)$ with $H^2(T_i,M)$.
Given three degree-one local classes, use (D.76.14)
on their cups with $u$, then use surjectivity in
(D.76.13). Restriction and cup commute, and local
periodicity gives exactly the original three classes.
Thus
\[
 H^1(E,M)\longrightarrow\bigoplus_iH^1(T_i,M)
       \text{ is onto},\qquad
 H^2(E,\{T_i\};M)=0.                                        \tag{D.76.15}
\]

\emph{Projectivity with compact module maps.}
For this specified augmentation kernel, the continuous
Hom construction and the relative long exact sequence
are the same ones used in D.74, now with $T_i$ in place
of the peripheral normalizers. The augmentation sequence
splits as compact $k$-vector spaces. Consequently
\[
 \operatorname{Ext}^1_R(\Delta_T,M)
   =H^1\bigl(E,\operatorname{Hom}_{k,\mathrm{cont}}(\Delta_T,M)\bigr)
   =H^2(E,\{T_i\};M)=0
                                                               \tag{D.76.16}
\]
for every finite $M$.
These are the compact/discrete Hom and relative
augmentation identifications of
\cite[Section~1 and Definition~2.1]{C-WilkesRelativeAccepted2018}.
They may also be read directly as follows. A compact
module extension of $\Delta_T$ by finite $M$ has a
continuous $k$-linear splitting; the difference of its
translates is a continuous cocycle with values in the
displayed Hom module. Vanishing of its class is
exactly an equivariant splitting.

For completeness, finite-coefficient vanishing here
does give projectivity in the compact module category.
Choose a surjection $p:R^m\to\Delta_T$, possible by
(D.76.10), and let $K=\ker p$.
For every open $R$-submodule $U$ of $R^m$ there is an
exact finite module sequence
\[
 0\longrightarrow (K+U)/U
 \longrightarrow R^m/U
 \longrightarrow\Delta_T/p(U)\longrightarrow0 .
                                                               \tag{D.76.17}
\]
Its pullback along $\Delta_T\to\Delta_T/p(U)$ splits
by (D.76.16). It gives a continuous $R$-linear map
$\ell_U:\Delta_T\to R^m/U$ lifting that quotient map.
The possible maps are a nonempty finite set, since
$\Delta_T$ is finitely generated and the target is finite.
Reduction for a finer $U$ carries such a lift to a
lift at every coarser $U$.
Any finitely many compatibility requirements are
realized at a common finer submodule. Compactness
of the product of these finite sets supplies compatible
$\ell_U$ for all $U$. Their inverse limit is an actual
continuous $R$-linear section of $p$.
Therefore $\Delta_T$ is a direct summand of $R^m$,
and is projective.

\emph{The full coefficient count and finite free quotients.}
Apply $\operatorname{Hom}_R(\Delta_T,-)$ to (D.76.11).
Projectivity gives an exact short sequence whose first
and last spaces are $\operatorname{Hom}_R(\Delta_T,M)$.
Frobenius reciprocity and (D.76.10) identify its middle
space with
$\operatorname{Hom}_{R_F}(R_F^4,M)$.
The reciprocity maps, using the coinduced functions in
(D.76.12), are evaluation at $1$ and its inverse
\[
        q\longmapsto
            \bigl[d\longmapsto(a\longmapsto q(ad))\bigr].
                                                               \tag{D.76.18}
\]
These are inverse continuous equivariant maps.
All three Hom spaces are finite. Their dimensions give
$2\dim\operatorname{Hom}_R(\Delta_T,M)=4\dim M$,
which is (D.76.3).

For an open normal subgroup $N\triangleleft E$, write
\[
 R_N=k[E/N],\qquad
 P_N=R_N\widehat\otimes_R\Delta_T
     =\Delta_T/\overline{\langle(n-1)d:n\in N,\ d\in\Delta_T\rangle}.
                                                               \tag{D.76.19}
\]
This is a finite projective left $R_N$-module:
base change of the direct-summand description above
is a direct summand of the finite free module $R_N^m$.
For every simple $R_N$-module $M$, (D.76.3) says
\[
       \dim_k\operatorname{Hom}_{R_N}(P_N,M)
         =2\dim_k M
         =\dim_k\operatorname{Hom}_{R_N}(R_N^2,M).
                                                               \tag{D.76.20}
\]
Let $J_N$ be the Jacobson radical of the finite algebra
$R_N$. Both Hom spaces factor through the respective
semisimple tops. For every simple $M$, its multiplicity
in such a top is the dimension of its Hom space divided
by $\dim_k\operatorname{End}_{R_N}(M)$.
Thus (D.76.20) identifies the two semisimple tops.
Lift that identification to a map $R_N^2\to P_N$.
Its cokernel has zero top, so is zero by Nakayama's
lemma, since $J_N$ is nilpotent.
Projectivity splits this surjection. The complementary
summand again has zero top and hence is zero.
This proves
$P_N\cong R_N^2$ as actual finite modules.

Let $\mathcal I_N=\operatorname{Iso}_{R_N}(R_N^2,P_N)$.
These are nonempty finite sets.
For $N'\subseteq N$, the canonical base-change identity
$R_N\otimes_{R_{N'}}P_{N'}=P_N$ sends an isomorphism
at $N'$ to one at $N$.
Every finite list of compatibility equations has a
solution obtained at a common finer open normal subgroup.
The finite intersection property therefore gives a
compatible family $(\phi_N)_N$.
The kernels in (D.76.19) form a neighbourhood basis
for $\Delta_T$: every finite module quotient has an
open normal subgroup acting trivially, and its kernel
contains the corresponding coinvariant kernel.
Consequently
\[
   \Delta_T=\varprojlim_NP_N,\qquad
   R^2=\varprojlim_NR_N^2,\qquad
   \phi=\varprojlim_N\phi_N
                                                               \tag{D.76.21}
\]
give the isomorphism (D.76.2), with inverse
$\varprojlim_N\phi_N^{-1}$.
This argument uses neither surjectivity of the maps
between the finite isomorphism sets nor an assumed
preferred basis.

\emph{The specified incidence map and its exact cokernel.}
Put $P_T=\bigoplus_i k[[E/T_i]]$.
The map
\[
        P_T\longrightarrow k[[E/H]],\qquad
                           gT_i\longmapsto gH
                                                               \tag{D.76.22}
\]
is well defined, continuous and equivariant.
It kills the image of $d_T$.
In $P_T/\operatorname{im}d_T$, the three identity
cosets are equal. Each $t_i$ fixes the common class,
because it fixes $e_i$ and $e_i=e_0$ there.
Continuity makes their closed generated group $H$
fix that class. Thus
$gH\mapsto[g e_0]$ is a well-defined continuous inverse
to the map induced by (D.76.22).
Checking the generators $gT_i$ and $gH$ proves both
compositions are the identity.
All images are closed, since the source modules are
compact. The augmentations commute with these maps.
Taking their kernels gives exactly the cokernel
isomorphism in (D.76.4).

The augmentation kernel of $k[[E/H]]$ vanishes exactly
when $E/H$ is a singleton. If two distinct cosets
exist, some finite clopen quotient separates them;
the difference of their two basis elements survives
in its finite permutation module.
Therefore $d_T$ is onto exactly when $H=E$.
Choose $\phi$ in (D.76.21) and set
$A_T=\phi^{-1}d_T:R^2\to R^2$.
If this map is onto, its reduction at every $N$
is an onto endomorphism of the finite module $R_N^2$,
hence is invertible. Taking the compatible inverse
limits makes $A_T$ invertible. The converse is immediate.
This proves the stated exact incidence criterion.
It supplies no reason to replace $d_T$ by the unrelated
isomorphism $\phi$.

\emph{The literal generators and the remaining removal target.}
The normalized action gives
$f\alpha(f)=1$ and
$\alpha(y)=f^{-1}y^{-1}f$.
Direct substitution shows that every element in
(D.76.5) squares to one.
They generate the original generators through
\[
            x=(xs)s,\qquad y=(yfs)(fs),\qquad s=s.
                                                               \tag{D.76.23}
\]
D.74's unique equation $h^{-1}\alpha(h)=f$ gives
$fs=h^{-1}sh$.
The other two involutions have respectively the
second and third section classes of D.34: replacing
$x^{e_2}s$ by $xs$, and $y^{e_2}fs$ by $yfs$, is
conjugation by $x^{(1-e_2)/2}$ and $y^{(1-e_2)/2}$.
These full half-exponents exist because $1-e_2$
has zero dyadic component.
Thus the four generators have exactly the duplicated
class stated in the claim.

Set $Y=hyh^{-1}$. The coboundary equation gives
$\alpha(Y)=Y^{-1}$, and
$h(yfs)h^{-1}=Ys$.
Hence the three elements in (D.76.6) represent the
three actual classes.
The subgroup $K=\overline{\langle x,Y\rangle}$ is
$\alpha$-invariant, and
\[
       \overline{\langle s,xs,Ys\rangle}
                          =K\langle s\rangle,\qquad
       K\cap\langle s\rangle=1,\qquad
       H\cap F=K.
                                                               \tag{D.76.24}
\]
These follow from $x=(xs)s$, $Y=(Ys)s$ and the
literal quotient label $u(s)=1$.
Equations (D.76.4) and (D.76.24) specify the exact
three-generator removal question.
Even its solution would require the original
third-cusp equation to give the peripheral conjugator
of D.35. Neither that equation nor a new group
presentation is inferred from the free module (D.76.2).
\end{proof}
